\subsection{From successful inputs to a common packet schedule}

Fix a packet size $s\ge1$. Use the hypothesis of \cref{thm:main} with $\eta=1/(2s)$.
We can choose $d$ as large as needed with $b\ge s$ and a nonempty success
set $G=G_{d,b}$ of deficit $h=L-\log_2|G|\le L/(2s)$. Write $b=as+r$
with $a\ge1$ and $0\le r<s$. Each matrix entry consists of $a$ complete
$s$-bit packets followed by $r$ tail bits. There are $n=d^2a$ complete
packets. The packet at $(i,j,\ell)$ is demanded at $(j,i,\ell)$ for $0\le\ell<a$. Since $L=sn+rd^2\le2sn$, we have $h\le n$.

Partition $G$ by the entire string of $rd^2$ tail bits. Some fixed
assignment $z$ leaves a set $F_z$ of successful packet arrays with
\[
  |F_z|\ge |G|\,2^{-rd^2}=2^{sn-h}.
\]
Thus fixing the tails preserves the bound $h$ on the absolute deficit,
with $sn$ as the new ambient bit count. Throughout this computation the
tail of an entry is the fixed string specified by its address in $z$.

The auxiliary input tape uses one cell per packet. The last cell of each
entry holds its $s$ payload bits followed by its fixed tail, so its length
in the original input is $s+r$. All other in-range cells have length $s$.
Each cell outside the input represents $s$ end markers, and each work cell
holds $s$ ordinary work symbols. Starting from any offset in a cell, $s$
source transitions can visit only that cell and its two neighbours.

One batch evaluates these transitions in eleven packet ticks. The input
and work heads move one cell left and read the left, centre and right
cells in succession. A stationary tick evaluates the batch on these
windows. The heads traverse the windows in reverse to write back the
three work cells, return to the old centres, and move by at most one cell
to the new centres. Other work writes preserve the scanned symbols. The final stationary tick may append one packet.
At the start of a batch the pending output is shorter than its next
cell. The batch adds at most $s$ raw output bits, so it completes at
most one cell. On completion, emit its first $s$ bits, pass over its
tail, and keep the suffix for the next cell. Every successful run emits
exactly the $n$ payload packets in transposed order. The public cell
lengths and tail values belong to this auxiliary computation.
\Cref{sec:simulation} will return to ordinary bit tapes.

Use $R=\lfloor\lfloor c(L+1)\rfloor/s\rfloor+1$ batches on every input,
keeping the source configuration fixed after halt. For
$J=2\lceil c\rceil+1$, the bounds $R\le J(n+1)$ and $T=11R\le11J(n+1)$
are uniform in $s,d,b$. Within each batch the varying part of the
schedule consists of the $k+1$ displacements to the new centres and the
packet-append flag. These give at most $(2\cdot3^{k+1})^R$ schedule
words. Hence for an integer $A_0=A_0(k,c)$ some fixed schedule has a
class $W$ satisfying
\[
  \log_2|W|\ \ge\ sn-h-A_0 n\ \ge\ sn-(A_0+1)n.
\]
The compiler further records the finite control, offsets, windows,
pending output, work writes and any output packet. They fit in
$\beta_s=O_M(s)$ bits per tick. On the fixed schedule the input address,
cell type, tick within the batch, and count of previously appended
packets are public. The record at a tick is determined by its input
packet, the preceding control record, and the records from the last
visits to the scanned work cells. Initial values and exterior input
symbols are fixed. This locality and the output at each designated
append tick are proved by \lean{TapeStorageRates.PaidClasses.compiler}.
